\section{Conclusions}
\label{sec:conclusions}

A test that counts the dates at which a behavioural parameter changed while a limit was breached turns a supervisory finding into a procedure with a known behaviour. On a parameter that a bank states and holds, it flags every change made at a breach date (\statedAltEngine{} of \statedN{} histories in the base cell, \statedAltMinLo{} at the lowest lower bound over \statedCells{} cells), it is silent when no change is made at a breach date, and it reproduces the stored rows of case C14 and agrees with its definition on \gridTotal{} enumerated and \randomN{} random histories. Its false-flag rate on unrelated changes is not a size to be chosen: it is $1-(1-\lambda)^m$ in expectation, from \statedNullMin{} to \statedNullMax{} over the cells measured, and the Fisher, Mann--Whitney and Welch tests, which hold a size below the nominal, detect only \opTestsMin{} to \opTestsMax{} of the planted changes at half of the breach dates at those lengths of history.

On a parameter that is re-estimated at each date, the exact test is red whenever a breach exists, up to the ties where the estimate sits at its cap: it flags \fitNullEngine{} of the null banks and its area against a planted override is \fitAucCEngine{}. The red rows of the stored case are therefore a statement about the estimator, as the product says, and they follow from Proposition~\ref{prop:degeneracy}. A tolerance on the size of a change, calibrated on a null population by a rank rule that holds its size on average in finite samples, restores the discrimination for overrides larger than the natural changes: \fitNullTol{} false flags and \fitAddDTol{} detection of a 4-year lengthening, with an area of \fitAucDTol{}. Calibrated at a breach prevalence of 0.25 and applied at 0.50 it flags \fitMisCTol{} of the null banks (\fitMisCTolLo{} to \fitMisCTolHi{}\,\%), a departure from 5\,\% that is inside the spread of a single calibration, so the size is held to within what a calibration can show; it is conservative on the real series; and an override toward the longest tranche is detected weakly on the simulated banks (no area above \wTopAuc{}, that of the break rule) and moderately on the real series (\wRealTopAuc{} for the level-shift test). The published test of a break at an unknown date, with the requirement that the break fall on a breach date, flags \fitAddCShiftB{} and \fitAddDShiftB{} of the 2- and 4-year lengthenings at \fitNullShiftB{} false flags, above 5\,\%, and \fitMisCShiftB{} (\fitMisCShiftBLo{} to \fitMisCShiftBHi{}\,\%) when it is calibrated at 0.25 and applied at 0.50, outside the spread of a calibration: the tolerance, at \fitNullTol{} false flags, ranks them better at both sizes and flags more of the 4-year ones. Three of the nine predictions failed: the conventional size on a stated parameter (H5), a ranking of a re-estimated parameter at the area predicted (H6) and, by its point estimate, the transfer of a tolerance across breach prevalences (H8).

The next steps follow from these limits. The first is to apply the exact test to the parameters that a bank states, including the overlay it adds to a fitted value, and to report its false-flag rate for the bank's own calendar from Equation~\eqref{eq:false-flag}. The second is to implement the tolerance rule with a calibration on a null population built from the bank's own estimator and breach process, and to test it on histories with a known override. The third is to extend the experiments to a longer panel of banks with their recorded limit frameworks, so that the rates are measured where the behaviour is real.
